\documentclass[10pt]{article}
\providecommand{\FIGDIR}{fig}
\newcommand{\DIGESTDATE}{5 September 2026}
% ===================== PM Research Watch - shared preamble =====================
\usepackage[T1]{fontenc}
\usepackage[utf8]{inputenc}
\usepackage{amsmath,amssymb}
\usepackage{tgpagella}
\usepackage{tgheros}
\renewcommand{\sfdefault}{qhv}
\renewcommand{\ttdefault}{lmtt}
\usepackage{microtype}
\usepackage[english]{babel}

\usepackage[a4paper,top=17mm,bottom=17mm,left=17mm,right=17mm,
            headsep=6mm,footskip=10mm]{geometry}
\usepackage{graphicx}
\usepackage{xcolor}
\usepackage{colortbl}
\usepackage{tikz}
\usetikzlibrary{calc}
\usepackage[most]{tcolorbox}
\usepackage{titlesec}
\usepackage{fancyhdr}
\usepackage{booktabs}
\usepackage{array}
\usepackage{longtable}
\usepackage{enumitem}
\usepackage{multicol}
\usepackage{ragged2e}
\usepackage[hidelinks]{hyperref}

% ------------------------------------------------------------------ palette
\definecolor{Ink}    {HTML}{14262E}
\definecolor{Deep}   {HTML}{0F4C5C}
\definecolor{Teal}   {HTML}{1F7A8C}
\definecolor{Sky}    {HTML}{4FA3C4}
\definecolor{Sage}   {HTML}{6FA287}
\definecolor{Amber}  {HTML}{E3A008}
\definecolor{Clay}   {HTML}{B5643C}
\definecolor{Coral}  {HTML}{D1495B}
\definecolor{Violet} {HTML}{6A5B9E}
\definecolor{Slate}  {HTML}{4A6572}
\definecolor{Paper}  {HTML}{FBF9F5}
\definecolor{Mist}   {HTML}{EEEAE2}
\definecolor{Rule}   {HTML}{D9D3C9}

\pagecolor{Paper}
\color{Ink}

\hypersetup{colorlinks=true, linkcolor=Deep, urlcolor=Teal, citecolor=Teal,
            pdftitle={PM Research Watch}, pdfauthor={Automated daily digest}}

% ------------------------------------------------------------------ headings
\titleformat{\section}
  {\sffamily\bfseries\large\color{Deep}}{}{0pt}
  {}[{\vspace{1.2mm}\color{Amber}\titlerule[1.1pt]}]
\titlespacing*{\section}{0pt}{5.5mm}{2.8mm}

\titleformat{\subsection}
  {\sffamily\bfseries\normalsize\color{Teal}}{}{0pt}{}
\titlespacing*{\subsection}{0pt}{3.6mm}{1.4mm}

\setlist[itemize]{leftmargin=4.2mm, itemsep=1.1mm, topsep=1.3mm, parsep=0pt,
                  label=\textcolor{Teal}{\footnotesize$\blacktriangleright$}}

% ------------------------------------------------------------------ header/footer
\pagestyle{fancy}
\fancyhf{}
\renewcommand{\headrulewidth}{0pt}
\renewcommand{\footrulewidth}{0pt}
\fancyhead[L]{\sffamily\scriptsize\color{Slate}PM RESEARCH WATCH\;\textcolor{Rule}{\textbar}\;\DIGESTDATE}
\fancyhead[R]{\sffamily\scriptsize\color{Slate}Particulate matter monitoring \& health effects}
\fancyfoot[C]{\sffamily\scriptsize\color{Slate}\thepage}

% ------------------------------------------------------------------ components
% metric tile
\newcommand{\metric}[3]{%
  \begin{tcolorbox}[enhanced, colback=white, colframe=Rule, boxrule=0.5pt,
    arc=1.6mm, left=2.2mm, right=2.2mm, top=1.8mm, bottom=1.8mm,
    borderline west={1.7mm}{0pt}{#3}, nobeforeafter, valign=center,
    height=17mm, width=\linewidth]
    \RaggedRight
    {\sffamily\bfseries\fontsize{17}{17}\selectfont\color{#3}#1}\\[0.5mm]
    {\sffamily\fontsize{6.5}{7.4}\selectfont\color{Slate}#2}
  \end{tcolorbox}}

% pull-quote / key insight box
\newtcolorbox{keybox}[1][]{enhanced, colback=Mist, colframe=Mist,
  arc=1.6mm, left=3.2mm, right=3.2mm, top=2.4mm, bottom=2.4mm,
  borderline west={1.7mm}{0pt}{Amber}, #1}

% subtopic band
\newcommand{\band}[2]{%
  \vspace{3.2mm}\noindent
  \begin{tcolorbox}[enhanced, colback=#2, colframe=#2, arc=1.2mm,
    left=2.6mm, right=2.6mm, top=1.1mm, bottom=1.1mm, boxrule=0pt,
    width=\linewidth, nobeforeafter]
    {\sffamily\bfseries\small\color{white}#1}
  \end{tcolorbox}\vspace{1.6mm}}

% one paper entry: {short cite}{title}{journal}{doi}{body}
\newcommand{\paperentry}[5]{%
  \par\vspace{1.9mm}\noindent
  {\sffamily\bfseries\footnotesize\color{Ink}#1}\;%
  {\sffamily\fontsize{7}{8}\selectfont\color{Slate}#3}\\[0.4mm]
  {\footnotesize\itshape\color{Deep}#2}\\[0.6mm]
  {\footnotesize #4}\\[0.3mm]
  {\fontsize{6.6}{7.6}\selectfont\color{Slate}\texttt{doi:}\,\href{https://doi.org/#5}{\texttt{\detokenize{#5}}}}
  \par}

% caption style
\newcommand{\figcap}[1]{\par\vspace{1.0mm}
  {\fontsize{7.2}{8.6}\selectfont\color{Slate}\textsf{\textbf{Fig.}}\; #1\par}\vspace{1.5mm}}

\providecommand{\permil}{\textperthousand}

\setlength{\parindent}{0pt}
\setlength{\parskip}{0pt}
\renewcommand{\arraystretch}{1.18}


\begin{document}

% ==================================================================== MASTHEAD
\thispagestyle{empty}
\begin{tikzpicture}[remember picture, overlay]
  \fill[Deep] (current page.north west) rectangle
      ($(current page.north east)+(0,-67mm)$);
  \fill[Amber] ($(current page.north west)+(0,-67mm)$) rectangle
      ($(current page.north east)+(0,-68.4mm)$);
  \foreach \x/\y/\r in {22/12/0.9, 41/26/1.5, 63/9/0.7, 88/31/1.1, 112/17/1.9,
                        138/29/0.8, 159/13/1.3, 178/54/1.0, 196/21/1.6,
                        30/44/0.6, 74/57/1.2, 125/48/0.9, 168/6/0.8, 191/59/0.7,
                        52/18/1.0, 100/60/0.6, 148/20/1.1, 66/49/1.4, 205/45/0.8}
    \fill[Sky, opacity=0.30] ($(current page.north west)+(\x mm,-\y mm)$) circle (\r mm);
\end{tikzpicture}

\vspace*{4mm}
{\sffamily\fontsize{7.6}{9}\selectfont\color{Sky}%
 AUTOMATED DAILY LITERATURE DIGEST\;\textbullet\;PhD RESEARCH WATCH\par}
\vspace{2.2mm}
{\sffamily\bfseries\fontsize{27}{29}\selectfont\color{white}%
 Particulate Matter\par}
{\sffamily\fontsize{27}{29}\selectfont\color{Sky}%
 Monitoring \& Health Effects\par}
\vspace{5.0mm}
{\sffamily\fontsize{8.4}{10}\selectfont\color{white}%
 Issue 2026-09-05\;\;\textcolor{Amber}{\textbar}\;\;Entry window 5 September 2026\;\;%
 \textcolor{Amber}{\textbar}\;\;7 records screened in scope\par}
\vspace{18mm}

% -------------------------------------------------------------- metric strip
\noindent
\begin{minipage}[t]{0.190\linewidth}\metric{7}{RECORDS IN SCOPE\\(22 EXCLUDED)}{Deep}\end{minipage}\hfill
\begin{minipage}[t]{0.190\linewidth}\metric{0}{PubMed [EDAT] HITS\\ON EITHER AXIS}{Coral}\end{minipage}\hfill
\begin{minipage}[t]{0.190\linewidth}\metric{5}{OF 7 ON THE SENSING /\\EXPOSURE AXIS}{Amber}\end{minipage}\hfill
\begin{minipage}[t]{0.190\linewidth}\metric{2}{WITH A HUMAN\\HEALTH ENDPOINT}{Teal}\end{minipage}\hfill
\begin{minipage}[t]{0.190\linewidth}\metric{2}{METADATA ONLY\\(NO ABSTRACT)}{Clay}\end{minipage}

\vspace{6mm}

% -------------------------------------------------------------- signal of day
\section*{\sffamily\bfseries\color{Deep}Signal of the day}
\vspace{-1.5mm}{\color{Amber}\rule{\linewidth}{1.1pt}}\vspace{2.5mm}

\begin{keybox}
{\sffamily\bfseries\small\color{Deep}Aviation and roadway leave separable fingerprints on
the same neighbourhood --- but only if the network measures particle \emph{number} and
absorption, not mass.}\\[1.4mm]
\footnotesize
Sawan et al.\ (\textit{Environ Sci Technol}) ran a hybrid mobile and fixed-site campaign in
a marginalized northwest Toronto neighbourhood sitting downwind of Toronto Pearson and a
freight corridor, and showed the two dominant combustion sources decouple spatially.
Ultrafine particle number tracked flight activity (\textbf{$r=0.58$, $p<0.05$ at 3.5\,km})
with \textbf{1.7--3.3-fold enrichment} under downwind conditions, still
\textbf{3.2-fold at residential sites $\sim$8\,km} from the airport; black carbon instead
followed roadway proximity and freight, correlating weakly with flights. The diagnostic
detail is that BC--UFP correlation collapses from $r=0.62$ under northerly winds to
$r=0.1$ downwind of the airport --- shared sources versus independent ones, visible only
because both pollutants were measured together. Random Forest models
($R^2=0.79$ UFP, $0.87$ BC) separated airport proximity as a UFP-specific driver from
freight trucking as a BC-specific one. For distributed low-cost networks the implication
is unwelcome and concrete: a PM$_{2.5}$-mass node sited here would have reported an
essentially null aviation signal, because the aviation contribution lives in a size range
that contributes negligible mass. Source separation of this kind needs a CPC-class number
channel and an absorption channel at the same node, which is precisely the pairing most
low-cost deployments omit.
\end{keybox}

\vspace{3mm}

% -------------------------------------------------------------- exec brief
\section*{\sffamily\bfseries\color{Deep}What changed today}
\vspace{-1.5mm}{\color{Amber}\rule{\linewidth}{1.1pt}}\vspace{2.5mm}

\footnotesize
\begin{itemize}
\item \textbf{PubMed returned a true zero, verified two ways.} No record entering PubMed on
5 September matched either the health axis or the sensing axis. This is not a harvester
fault: 544 records entered PubMed that day against 4{,}765 on Friday 4 September, and the
independent PubMed connector run with looser field tags also returned 0/0. Saturday
deposits are thin and, this week, contained no PM work at all.

\item \textbf{The issue was carried entirely by the non-PubMed legs.} Crossref-by-ISSN
supplied four of seven records and Europe PMC three. Europe PMC returned 11 hits but
\textbf{seven were already shipped in the 4 September issue} --- its
\texttt{CREATION\_DATE} runs about a day behind PubMed \texttt{[EDAT]}, so cross-source
deduplication, not recall, is what that leg mostly exercises.

\item \textbf{Instrumentation dominates a health-thin day.} Five of seven records sit on
the sensing or exposure-assessment axis (71\%, second only to 20 August's 78\%), and only
\textbf{two records carry a human health endpoint at all}. No record in the window reports
an effect estimate with a confidence interval, so there is no forest plot today --- the
seventh such issue of 42.

\item \textbf{Two source-attribution results worth reading together.} Sawan et al.\ separate
aviation UFP from roadway BC by wind sector; Shi S et al.\ show that vanadium, the
canonical ship-plume tracer, is carried by only \textbf{35.8\% of plume particles} and that
coverage moves with vessel speed (29.5\% at 4.9 knots vs.\ 39.7\% at 13.4). Both are
warnings that a single-species or single-metric tracer under-identifies its own source.

\item \textbf{A speed-reduction policy may trade CO$_2$ for composition.} In Shi S et al.'s
matched-vessel pair the CO$_2$ benefit of slowing coincided with a carbon-rich particle
shift and lower detectable-vanadium coverage. They also find \textbf{$\sim$20\% of analysed
particles meet single-particle biomass-burning criteria}, i.e.\ ship plumes generate false
positives in exactly the source-apportionment schemes used to police them.

\item \textbf{Trial watch is a structural zero, not a topical one.} ClinicalTrials.gov
posted \emph{no} updates registry-wide on 5 September --- nor on 29/30 August --- against
973--1{,}001 on each weekday. Every Saturday and Sunday issue will carry this zero for
calendar reasons, and it should not be read as an absence of PM trial activity.

\item \textbf{Two records ship without abstracts.} The \textit{Atmos Pollut Res} tire-wear
tracer review and the \textit{Atmos Res} dust-aging lidar paper are indexed by Crossref
with no abstract and no Europe PMC record; both are titled informatively enough to
place but not to assess, and both go on the retry list.
\end{itemize}

\clearpage

% ==================================================================== ANALYTICS
\section*{\sffamily\bfseries\color{Deep}Corpus analytics}
\vspace{-1.5mm}{\color{Amber}\rule{\linewidth}{1.1pt}}\vspace{3mm}

\noindent\includegraphics[width=\linewidth]{\FIGDIR/f1_subtopics.png}
\figcap{Subtopic assignment of the seven in-scope records. Four of the ten clusters are
occupied --- tying the fewest recorded (5, 23 and 24 August) --- and at seven records this
is the \textbf{smallest single-day issue this watch has shipped}, below the previous
minimum of eight on 29 August. \textit{Sensing, forecasting \& instrumentation} leads with
three; six of the ten clusters, including all four of the cardiovascular, neuro,
respiratory and reproductive health blocks, are empty.}

\vspace{2mm}
\noindent\includegraphics[width=\linewidth]{\FIGDIR/f2_design_endpoint.png}
\figcap{Left --- study architecture. \textit{Measurement campaign} accounts for three of
seven, with one laboratory-combustion study, one in-vitro toxicology preprint and two
records that could not be classified beyond \textit{Metadata only} because no abstract has
been deposited. Right --- health endpoint. Five records carry no human endpoint at all
(two aerosol chemistry, one instrument, one exposure, one with no health endpoint), which
is why the endpoint panel is so sparse; only the Johannesburg occupational study and the
in-vitro preprint sit on a health axis.}

\vspace{2mm}
\noindent\includegraphics[width=\linewidth]{\FIGDIR/f3_metric_geography.png}
\figcap{Left --- particle metric. Four metrics span seven records: ultrafine number (2),
unspeciated PM and emission factors (3), lidar optical properties (1) and PM$_{2.5}$ (1).
\textbf{PM$_{2.5}$ mass appears alone exactly once, and only as a dosed in-vitro
concentration} (25--100\,$\mu$g/mL) rather than an ambient measurement --- every ambient
record here is instrumented on number, emission factors or optical properties instead.
Right --- geographic provenance. North America and
China contribute two each; the Johannesburg occupational study is the day's only
sub-Saharan African record, and South Asia, Latin America and Oceania are absent.}

\vspace{2mm}
\noindent\includegraphics[width=\linewidth]{\FIGDIR/f6_lifecourse.png}
\figcap{Life-course windows addressed. Only the working-age stage is occupied, by the
12-teacher Johannesburg personal-sampling study; the preconception, infancy, adolescence
and 65+ windows are all empty. A single-stage day is a direct consequence of a corpus
that is 71\% instrumentation and exposure assessment.}

\vspace{2mm}
\noindent\includegraphics[width=\linewidth]{\FIGDIR/f4_heatmap.png}
\figcap{Subtopic $\times$ study-architecture cross-tabulation. With seven records the
matrix is necessarily sparse, but the standing structural finding survives at small
sample: the \textit{Sensing} and \textit{Exposure} rows again contain \emph{no}
health-endpoint linkage, and the one mechanistic record is again purely experimental with
no bridge to human biomarker data. There is no forest plot in this issue --- no record in
the window reports an effect estimate with a confidence interval.}

% ==================================================================== DIGEST
\section*{\sffamily\bfseries\color{Deep}Paper-level digest}
\vspace{-1.5mm}{\color{Amber}\rule{\linewidth}{1.1pt}}\vspace{1mm}

{\fontsize{7.4}{9}\selectfont\color{Slate}Ordered by cluster. Each entry gives the
methodological core and the finding that matters, not an abstract paraphrase.\par}

\band{Sensing, forecasting \& instrumentation\hfill 3 records}{Amber}

\paperentry{Shi S et al.\ 2026}
{Vessel speed reduction trades CO$_2$ against plume composition --- and against the tracer
that identifies the plume}
{npj Clim Atmos Sci}
{UAV interception of 14 fresh ship plumes with single-particle mass spectrometry and
population-balance inversion to reconstruct ultrafine and coarse modes. Two marker gaps
emerge. \textbf{Only 35.8\% of the reconstructed particle-number population carried
detectable vanadium} (29.5\% at 4.9 knots vs.\ 39.7\% at 13.4 knots in a matched-vessel
pair; 0--70.8\% across a speed-blind manifold), so the canonical ship tracer is
operation-dependent rather than conserved. Separately, \textbf{$\sim$20\% of analysed
particles met biomass-burning criteria}, a false-positive channel in ship-influenced
samples. Sulfate-related ion signals ran 2.3--10.7-fold higher in V/VO-bearing particles.
The authors are careful that laser desorption/ionization matrix effects may contribute to
that contrast, which is the right caveat and limits how far the V--S co-location can be
pushed mechanistically. Single campaign, 14 plumes, one matched pair --- the speed
comparison rests on very few vessels.}
{10.1038/s41612-026-01532-3}

\paperentry{Wang S et al.\ 2026}
{Smoldering chaparral emits 10--70$\times$ the particulate of flaming, and fuel moisture
sets the multiplier}
{Fire Saf J}
{FTIR quantification of gas and particulate emissions from controlled flaming and
smoldering combustion of four dominant California chaparral species at varying fuel
moisture. \textbf{Smoldering produced 10--70-fold higher PM emission factors than flaming
across all species}, with the amplification ratio itself varying by fuel type and
moisture; higher moisture enhanced incomplete-combustion products, principally CO and
hydrocarbons, with live eucalyptus highest at
\textbf{37.18 $\pm$ 0.65\,g/kg} hydrocarbons under smoldering. Cross-ecosystem comparison
found chaparral broadly similar to western shrubland and grassland fuels except in
nitrogen-containing species. The limitation that matters is scale: these are
bench-scale burns, and the authors note combustion scale and mixing conditions shift
emission profiles, so transfer to landscape fire inventories is an assumption rather than
a result. Directly useful for anyone weighting wildfire smoke exposure by combustion
phase rather than by total burned area.}
{10.1016/j.firesaf.2026.104687}

\paperentry{Ji et al.\ 2026}
{Dust aerosol aging tracked by combined Mie, polarization and fluorescence lidar}
{Atmos Res}
{\textbf{Metadata only.} Crossref carries the title, author list (8) and 5 September
registration date but no abstract, and there is no Europe PMC record; the publisher's
full text was not retrievable on this run. From the title the study links optical
evolution during dust aging to chemical transformation using a three-channel lidar ---
a combination relevant to anyone attempting composition inference from optical
instruments rather than mass. \textbf{No quantities, site, campaign length or retrieval
performance are asserted here, because none could be verified.} Added to the abstract
retry list.}
{10.1016/j.atmosres.2026.109312}

\band{Exposure assessment \& modelling\hfill 2 records}{Amber}

\paperentry{Sawan et al.\ 2026}
{Aviation drives ultrafine number, freight drives black carbon, in the same
neighbourhood}
{Environ Sci Technol}
{Hybrid mobile and fixed-site campaign in a marginalized northwest Toronto neighbourhood
downwind of Toronto Pearson and a freight corridor. UFP correlated with flight activity
(\textbf{$r=0.58$, $p<0.05$ at 3.5\,km}) and was enriched \textbf{1.7--3.3-fold} downwind
of the airport, measured as the ratio of interquartile ranges under downwind versus other
wind conditions, reaching \textbf{3.2-fold at residential sites $\sim$8\,km away}. BC
instead tracked roadway proximity and freight and correlated weakly with flight counts.
BC--UFP correlation was strong under northerly winds (mean $r=0.62$) but collapsed to
$r=0.1$ downwind of the airport. Random Forest models ($R^2=0.79$ UFP, $0.87$ BC)
identified industrial land use and a highway as shared drivers, airport proximity as
UFP-specific and freight trucking as BC-specific. The inference is observational and
wind-sector stratified, so airport attribution rests on a co-location argument rather than
on any tracer unique to jet exhaust. The lesson for network design is that mass-based
siting would have missed the aviation signal entirely.}
{10.1021/acs.est.6c06666}

\paperentry{Scalabrin et al.\ 2026}
{Chemical tracers for tire wear particles in atmospheric aerosol: analytical approaches
and challenges}
{Atmos Pollut Res}
{\textbf{Metadata only.} Crossref registers the title, 11 authors and a 5 September
creation date with no abstract, and the record has no Europe PMC counterpart. The title
places it as a review of tracer chemistry for apportioning non-exhaust tire wear ---
a live problem, since tire wear is a growing share of traffic PM precisely as exhaust
controls tighten, and the standard markers each fail differently. \textbf{No findings,
recommended tracer or study count are asserted here.} A targeted Scholar Gateway
full-text query returned only adjacent June 2026 reviews, not this paper. Added to the
abstract retry list.}
{10.1016/j.apr.2026.103194}

\band{Occupational \& indoor exposure\hfill 1 record}{Clay}

\paperentry{Mbazima et al.\ 2026}
{Chalk dust puts teachers over a deposition threshold while staying under the
occupational limit}
{Int J Environ Health Res}
{Sixty personal breathing-zone samples from 12 Johannesburg secondary-school teachers
(7 female, 5 male), collected Monday--Friday 08:00--14:35 with GilAir-3 pumps, a 4\,$\mu$m
cut-point cyclone and 25\,mm PVC filters at 1.85\,L/min, then run through a multiple-path
particle dosimetry model at low, moderate and high exertion for nasal and oronasal
breathing. \textbf{PM$_4$ ranged 0.019--0.761\,mg/m$^3$, mean 0.126 $\pm$
0.134\,mg/m$^3$}; chalk use ($\beta=0.468$) and sampling duration ($\beta=0.305$, both
$p<0.001$) predicted concentration. Breathing route dominated deposition, and alveolar
clearance was slower than tracheobronchial in both sexes. The authors' conclusion --- that
risk persists below the occupational exposure limit --- rests on modelled rather than
measured dose, and 12 teachers in one region is a small and non-random base. Note the
standard deviation exceeds the mean, so the distribution is heavily right-skewed and the
mean is a poor summary.}
{10.1080/09603123.2026.2726420}

\band{Mechanistic toxicology\hfill 1 record}{Violet}

\paperentry{Han G et al.\ 2026 (preprint)}
{Picropodophyllotoxin attenuates PM$_{2.5}$ injury in pulmonary endothelial cells via
mTOR-dependent autophagy}
{Research Square (preprint)}
{Human pulmonary artery endothelial cells exposed to PM$_{2.5}$ at 25--100\,$\mu$g/mL for
24--48\,h with picropodophyllotoxin (2--50\,$\mu$M) or dexamethasone (50\,$\mu$M).
PM$_{2.5}$ reduced viability dose- and time-dependently with increased LDH leakage and ROS
and suppressed SOD/CAT and SGK1; PPT restored viability, SGK1 and mTOR phosphorylation and
attenuated TLR4, MyD88, LC3II and Beclin 1 upregulation. \textbf{Not peer reviewed, single
cell line, no animal or human data, and no dose--response statistics reported in the
abstract.} The exposure concentrations are far above realistic alveolar deposition, which
is the standard weakness of this design and limits the finding to pathway plausibility
rather than therapeutic inference. Included for completeness of the mechanistic map,
not as evidence of effect.}
{10.21203/rs.3.rs-10446885/v1}

\clearpage

% ==================================================================== TABLE
\section*{\sffamily\bfseries\color{Deep}Corpus register}
\vspace{-1.5mm}{\color{Amber}\rule{\linewidth}{1.1pt}}\vspace{2.5mm}

{\fontsize{6.7}{8.4}\selectfont
\setlength{\tabcolsep}{2.4pt}
\begin{longtable}{@{}>{\RaggedRight}p{25mm}>{\RaggedRight}p{41mm}>{\RaggedRight}p{27mm}%
>{\RaggedRight}p{23mm}>{\RaggedRight}p{18mm}>{\RaggedRight}p{15mm}@{}}
\toprule
\textsf{\textbf{First author}} & \textsf{\textbf{Focus}} & \textsf{\textbf{Journal}} &
\textsf{\textbf{Design}} & \textsf{\textbf{Metric}} & \textsf{\textbf{Region}}\\
\midrule
\endfirsthead
\toprule
\textsf{\textbf{First author}} & \textsf{\textbf{Focus}} & \textsf{\textbf{Journal}} &
\textsf{\textbf{Design}} & \textsf{\textbf{Metric}} & \textsf{\textbf{Region}}\\
\midrule
\endhead
\multicolumn{6}{@{}l@{}}{\cellcolor{Mist}\textsf{\textbf{\textcolor{Deep}{Neuro / mental health}}}}\\[0.6mm]
\href{https://doi.org/10.1001/jamanetworkopen.2026.25015}{Dales} & Wildfire-sourced PM$_{2.5}$ and ED visits for migraine and other primary headache & \textit{JAMA Netw Open} & Case-crossover & PM$_{2.5}$ (wildfire) & Canada\\
\href{https://doi.org/10.1093/ageing/afag215}{Glover} & Umbrella meta-review of urban environmental exposures and cognitive health & \textit{Age Ageing} & Umbrella review / meta-analysis & PM$_{2.5}$ & Global\\
\href{https://doi.org/10.1016/j.nsa.2026.107019}{Goossens} & In utero air pollution exposure and pre-/postnatal brain development & \textit{Neurosci Appl} & Narrative review & PM$_{2.5}$ / PM$_{10}$ / PM$_{\mathrm{coarse}}$ & Global\\
\href{https://doi.org/10.1016/j.isci.2026.116617}{Ma} & PM$_{2.5}$-induced neurotoxicity and depression: exposure risks and mechanisms & \textit{iScience} & Narrative review & PM$_{2.5}$ + components & Global\\
\href{https://doi.org/10.3760/cma.j.cn112150-20250804-00752}{Pan} & Maternal PM exposure and offspring neurodevelopment: effects and mechanisms & \textit{Zhonghua Yu Fang Yi Xue Za Zhi} & Narrative review & PM (bulk) & China\\
\addlinespace[0.9mm]
\multicolumn{6}{@{}l@{}}{\cellcolor{Mist}\textsf{\textbf{\textcolor{Teal}{Cardiovascular \& metabolic}}}}\\[0.6mm]
\href{https://doi.org/10.3760/cma.j.cn112338-20251120-00843}{Huang} & PM$_{2.5}$ and dyslipidaemia in rural Chinese children; dietary effect modification & \textit{Zhonghua Liu Xing Bing Xue Za Zhi} & Cross-sectional (multistage) & PM$_{2.5}$ & China\\
\href{https://doi.org/10.32394/pe/218966}{Kazmierska} & Air pollution and cardiovascular disease: review of current data & \textit{Przegl Epidemiol} & Narrative review & PM$_{2.5}$ / PM$_{10}$ & Global\\
\href{https://doi.org/10.1007/s00125-026-06814-2}{Xi} & Air pollution, built environment and temperature vs. prediabetes phenotypes & \textit{Diabetologia} & Prospective cohort & PM$_{2.5}$ & Germany\\
\href{https://doi.org/10.1161/JAHA.125.046917}{Yang} & Short-term PM$_{2.5}$ and hospitalisation/ED use in atrial fibrillation & \textit{J Am Heart Assoc} & Case-crossover & PM$_{2.5}$ (wildfire-influenced) & USA\\
\addlinespace[0.9mm]
\multicolumn{6}{@{}l@{}}{\cellcolor{Mist}\textsf{\textbf{\textcolor{Sky}{Reproductive \& developmental}}}}\\[0.6mm]
\href{https://doi.org/10.1016/j.jes.2025.10.043}{Chen} & Prenatal air pollution, fetal growth trajectory phenotypes and child neurodevelopment & \textit{J Environ Sci (China)} & Prospective cohort & PM$_{2.5}$ / PM$_{10}$ & China\\
\href{https://doi.org/10.1016/j.ecoenv.2026.120566}{Huang} & Joint prenatal bisphenol + air-pollutant mixture exposure and preschool IQ & \textit{Ecotoxicol Environ Saf} & Prospective cohort (mixtures) & PM$_{2.5}$ / PM$_{10}$ & China\\
\href{https://doi.org/10.1007/s40572-026-00552-8}{Peterson} & Ambient air pollution, PAHs and pubertal development: critical review & \textit{Curr Environ Health Rep} & Critical review & PM$_{2.5}$ components & Global\\
\href{https://doi.org/10.1016/j.envres.2026.125313}{Xie} & Long-term PM$_{2.5}$ and timing of pubertal development in boys and girls & \textit{Environ Res} & Prospective cohort & PM$_{2.5}$ & China\\
\addlinespace[0.9mm]
\multicolumn{6}{@{}l@{}}{\cellcolor{Mist}\textsf{\textbf{\textcolor{Sage}{Respiratory \& allergic}}}}\\[0.6mm]
\href{https://doi.org/10.1007/s00484-026-03289-z}{Bayar Muluk} & Climate change and otolaryngology: airway disease impacts and emission reduction & \textit{Int J Biometeorol} & Narrative review & PM$_{2.5}$ & Global\\
\href{https://doi.org/10.4103/mgr.MEDGASRES-D-25-00164}{Liao} & Short-term air pollution, respiratory symptoms and lung function in asthma & \textit{Med Gas Res} & Panel study & PM$_{2.5}$ / PM$_{10}$ & China\\
\href{https://doi.org/10.3390/epidemiologia7040097}{Sanchez de Prada} & Ambient PM and asthma-coded admissions in children 0-3 y, Spain & \textit{Epidemiologia} & Ecological panel & PM$_{2.5}$ / PM$_{10}$ & Spain\\
\href{https://doi.org/10.1016/j.jhazmat.2026.143076}{Yang} & Allergenic vs. toxic risk of airborne seed hairs compared with ambient PM$_{2.5}$ & \textit{J Hazard Mater} & Field measurement & PM$_{2.5}$ + bioaerosol & China\\
\addlinespace[0.9mm]
\multicolumn{6}{@{}l@{}}{\cellcolor{Mist}\textsf{\textbf{\textcolor{Violet}{Mechanistic toxicology}}}}\\[0.6mm]
\href{https://doi.org/10.1016/j.jes.2025.12.074}{Ding} & PM$_{2.5}$ promotes vascular calcification via miR-27b-5p/HES1 signalling & \textit{J Environ Sci (China)} & In vitro + organoid + murine & PM$_{2.5}$ & China\\
\href{https://doi.org/10.1007/s00360-026-01698-5}{Maraschi} & Osmoregulatory disruption and metal bioaccumulation from settleable atmospheric PM & \textit{J Comp Physiol B} & Controlled exposure (ecotox) & Settleable APM & Brazil\\
\href{https://doi.org/10.1016/j.jes.2025.11.018}{Wu} & Macrophage NLRP3 inflammasome in real-ambient PM-induced hepatic glucose dysregulation & \textit{J Environ Sci (China)} & Real-ambient murine + co-culture & Real-ambient PM & China\\
\addlinespace[0.9mm]
\multicolumn{6}{@{}l@{}}{\cellcolor{Mist}\textsf{\textbf{\textcolor{Clay}{Occupational \& indoor}}}}\\[0.6mm]
\href{https://doi.org/10.3389/fmedt.2026.1796828}{Chen} & Nested suction instrument to cut intraoperative surgical-smoke PM$_{2.5}$ & \textit{Front Med Technol} & Bench / device experiment & PM$_{2.5}$ (surgical smoke) & China\\
\href{https://doi.org/10.1093/annweh/wxag055}{Cole} & Determinants of respirable crystalline silica exposure in Australian tunnelling & \textit{Ann Work Expo Health} & Observational exposure study & Respirable dust / RCS & Australia\\
\href{https://doi.org/10.1016/j.jhazmat.2026.143081}{Ding} & Stage-resolved hair Zn-isotope dynamics as a firefighting combustion-exposure tracer & \textit{J Hazard Mater} & Repeated-measures biomarker & Combustion-derived PM & China\\
\href{https://doi.org/10.3390/jox16040119}{Ramsden} & Xenobiotic hazards in aircraft cabin air & \textit{J Xenobiot} & Review + risk economics & Ultrafine particles & Global\\
\addlinespace[0.9mm]
\multicolumn{6}{@{}l@{}}{\cellcolor{Mist}\textsf{\textbf{\textcolor{Amber}{Sensing, forecasting \& instrumentation}}}}\\[0.6mm]
\href{https://doi.org/10.1016/j.jenvman.2026.130515}{Shi} & Adaptive spatiotemporal graph transformer for multi-site multi-step PM$_{2.5}$ forecasting & \textit{J Environ Manage} & Deep-learning model & PM$_{2.5}$ & China / India\\
\href{https://doi.org/10.3390/bios16070380}{Tsuchiya} & Seasonal and indoor-outdoor characteristics of fluorescent aerosol particles in offices & \textit{Biosensors} & Real-time sensor campaign & Bioaerosol / coarse & Japan\\
\addlinespace[0.9mm]
\multicolumn{6}{@{}l@{}}{\cellcolor{Mist}\textsf{\textbf{\textcolor{Amber}{Exposure assessment \& modelling}}}}\\[0.6mm]
\href{https://doi.org/10.1029/2025GH001621}{Solorzano-Araque} & Multiscale exposure modelling: LOTOS-EUROS + mobility surveys + particle toxicity & \textit{GeoHealth} & CTM + mobility integration & PM$_{2.5}$ / PM$_{10}$ & Colombia\\
\addlinespace[0.9mm]
\multicolumn{6}{@{}l@{}}{\cellcolor{Mist}\textsf{\textbf{\textcolor{Coral}{Burden, policy \& mitigation}}}}\\[0.6mm]
\href{https://doi.org/10.1002/hsr2.72886}{Karamzad} & Ischaemic heart disease burden in MENA 1990-2021; PM attributable fraction 28.2\% & \textit{Health Sci Rep} & GBD secondary analysis & Ambient PM & MENA\\
\href{https://doi.org/10.1371/journal.pone.0347989}{Rees} & HEPA filtration to reduce respiratory infection in care homes & \textit{PLoS One} & RCT process evaluation & Filtered PM & UK\\
\href{https://doi.org/10.1016/j.jes.2026.01.009}{Shan} & Policy-driven emission reductions and post-pandemic rebound, high-resolution inventories & \textit{J Environ Sci (China)} & Emission inventory & PM emissions & China\\
\addlinespace[0.9mm]
\multicolumn{6}{@{}l@{}}{\cellcolor{Mist}\textsf{\textbf{\textcolor{Slate}{Other clinical endpoints}}}}\\[0.6mm]
\href{https://doi.org/10.3390/dj14070440}{Lupita} & PM$_{2.5}$, preventive-care capacity and caries/edentulism burden across the EU & \textit{Dent J} & Ecological panel & PM$_{2.5}$ & EU-27\\
\href{https://doi.org/10.1016/j.jdin.2026.06.008}{Park} & Long-term PM exposure and risk of androgenetic alopecia & \textit{JAAD Int} & Nationwide cohort & PM$_{2.5}$ / PM$_{10}$ & South Korea\\
\href{https://doi.org/10.3171/2026.3.PEDS2615}{Sun} & Environmental exposure and area deprivation vs. paediatric hydrocephalus outcomes & \textit{J Neurosurg Pediatr} & Registry cohort + geospatial & PM$_{2.5}$ (satellite-derived) & USA\\
\addlinespace[0.9mm]
\bottomrule
\end{longtable}}


\vspace{4mm}
\begin{keybox}
{\sffamily\bfseries\footnotesize\color{Deep}Provenance and method}\\[1.2mm]
{\fontsize{7.2}{8.8}\selectfont
Entry window \textbf{5 September 2026}, a single day. \textbf{PubMed} (NCBI E-utilities)
was queried on both the health and sensing axes against \texttt{[Date - Entry]} and
returned \textbf{zero hits on each}; the PubMed connector, run independently with looser
field tags, also returned zero, and a control query confirmed only 544 records of any
kind entered PubMed that day against 4{,}765 the previous day. The issue therefore rests
on \textbf{Crossref-by-ISSN} (18 raw hits, 4 in scope) and \textbf{Europe PMC}
(\texttt{CREATION\_DATE}, 11 raw hits, 3 in scope after seven were found already published
in the 4 September issue). \textbf{arXiv} returned HTTP 200 but its newest submission
remains 2026-08-20 for the third consecutive run, so it carried nothing.
\textbf{OpenAlex} returned HTTP 429 (unchanged since 28 July). \textbf{Consensus} returned
three hits, all verified out of window by Crossref \texttt{created} date (2018, 2019,
2025) --- a nil for the seventh consecutive run.
\textbf{ClinicalTrials.gov} posted no updates registry-wide on 5 September; the same holds
for 29 and 30 August against 973--1{,}001 on each intervening weekday, so the weekend
trial zero is a posting-calendar artefact and is reported as such.
Twenty-two records were screened out with a logged reason: seven as duplicates of the
4 September issue and fifteen as off-topic or non-research (aquatic microplastics and
PFAS, a corrigendum, a correspondence letter, indoor lighting, urban heat island and
precipitation forecasting).
Subtopic, design, particle-metric and geography labels are assigned from title, abstract
and keywords; assignment is single-label by dominant contribution and is a judgement, not
a controlled vocabulary. Two records carry no abstract from any source and are labelled
\textit{Metadata only}; nothing is asserted about their findings.
All DOI links resolve to the publisher of record. Abstracts remain the copyright of their
respective publishers.\par}
\end{keybox}

\end{document}
